% ------------------------------------------------------------------ 4
\section{Three greedy strategies}\label{sec:greedy}

The three strategies differ in the greedy \emph{selection criterion}, that is, in which object they commit to at each step:
\begin{description}[leftmargin=1.2cm,style=nextline]
\item[G1] the cheapest shortest path between two terminal components (an MST of the distance network \(D\));
\item[G2] the cheapest shortest path from the single growing tree to a terminal outside it;
\item[G3] the triple of terminals whose connection through one Steiner vertex saves the most against the current
spanning tree.
\end{description}
G1 and G2 only ever join terminals in pairs. G3 is the only one that can choose a Steiner vertex because it serves
three terminals at once, and \cref{sec:gaps} shows that this is the step that pays off. All three end with the same
clean-up, which we prove safe first.

\begin{lemma}[clean-up]\label{lem:cleanup}
Let \(H\subseteq G\) be connected with \(T\subseteq V(H)\). Replacing \(H\) by a minimum spanning tree of
\(G[V(H)]\) and then deleting non-terminal leaves until none is left produces a Steiner tree of cost at most \(w(H)\).
\end{lemma}
\begin{proof}
\(H\) contains a spanning tree of \(V(H)\) of cost at most \(w(H)\) (weights are non-negative), and \(H\subseteq G[V(H)]\), so
the MST of \(G[V(H)]\) costs at most \(w(H)\). Deleting a non-terminal leaf keeps a tree that still contains \(T\) and removes
an edge of weight \(\ge0\). The clean-up runs Prim on the heap (\(O(m+n\log n)\)) and a leaf queue (\(O(n)\)).
\end{proof}

\subsection{G1: distance-network MST (Kou, Markowsky and Berman)}
\begin{algorithm}[H]
\For{\(i=1\) \KwTo \(k\)}{ Dijkstra from \(t_i\) (Fibonacci heap); store \(d(t_i,\cdot)\) and the predecessor tree\; }
build \(D\) with \(D_{ij}=d(t_i,t_j)\); \(M\gets\) Prim on \(D\) (Fibonacci heap)\;
\(H\gets\) union of the shortest paths of \(G\) for the edges of \(M\)\;
\Return \textsc{Clean-Up}\((H)\)\;
\caption{G1 (\texttt{src/greedy1\_kmb.cpp}).}
\end{algorithm}

\paragraph{Running time.} \(k\) Dijkstra runs cost \(O(k(m+n\log n))\) by \cref{sec:amortized}; Prim on \(D\) costs \(O(k^2)\);
expanding \(k-1\) paths costs \(O(kn)\); the clean-up costs \(O(m+n\log n)\). Total \(O(k(m+n\log n))\), polynomial.
The original paper states \(O(kn^2)\) because it uses an array-based Dijkstra.

\begin{theorem}\label{thm:kmb}
Let \(\ell\ge2\) be the number of leaves of an optimal Steiner tree \(S^\ast\). Then
\(w(\mathrm{MST}(D))\le\big(2-\tfrac{2}{\ell}\big)\OPT\), and G1 returns a tree of cost at most \(\big(2-\tfrac{2}{\ell}\big)\OPT\le\big(2-\tfrac2k\big)\OPT\).
\end{theorem}
\begin{proof}
Double every edge of \(S^\ast\). The result is Eulerian; an Euler tour \(W\) has cost \(2\,\OPT\) and visits the leaves
of \(S^\ast\) in some cyclic order \(L_1,\dots,L_\ell\). Cut \(W\) at the leaves into \(\ell\) walks \(W_j\) from \(L_j\) to
\(L_{j+1}\). The longest of them costs at least \(2\OPT/\ell\); delete it. The other \(\ell-1\) walks form one walk of cost at
most \(2\OPT(1-1/\ell)\) that still visits every terminal: a leaf is an end of two walks, and an inner vertex of degree
\(\delta\ge2\) appears in \(\delta\) walks, so deleting one walk removes neither. Listing the terminals in the order this walk
first reaches them gives a path in \(D\) whose edges are no longer than the walk segments between them (\(d\) is a shortest-path
metric). That path is a spanning tree of \(D\), so \(w(\mathrm{MST}(D))\le(2-2/\ell)\OPT\). The expanded subgraph \(H\) costs at most
\(w(\mathrm{MST}(D))\) (shared edges are counted once), and \cref{lem:cleanup} does not increase it. Since \(\ell\le k\), the bound
\(2-2/k\) follows.
\end{proof}

\subsection{G2: shortest-path heuristic (Takahashi and Matsuyama)}
\begin{algorithm}[H]
\(S\gets\{r\}\)\;
\While{\(T\not\subseteq V(S)\)}{
  multi-source Dijkstra with label 0 on every vertex of \(S\), stopped when the first terminal \(t\notin V(S)\) is extracted\;
  add the shortest path from \(S\) to \(t\) to \(S\)\;
}
\Return \textsc{Clean-Up}\((S)\)\;
\caption{G2 for one root \(r\) (\texttt{src/greedy2\_sph.cpp}). The solver runs every terminal as the root and keeps the best tree.}
\end{algorithm}

\paragraph{Running time.} Each of the \(k-1\) rounds is one Dijkstra, \(O(m+n\log n)\), and adding a path is \(O(n)\). One root
costs \(O(k(m+n\log n))\); all \(k\) roots cost \(O(k^2(m+n\log n))\). Stopping at the first terminal is safe because a Dijkstra
label is final once it is extracted, and it shortens the rounds in practice.

\begin{theorem}\label{thm:sph}
For every root, G2 returns a tree of cost at most \(w(\mathrm{MST}(D))\le\big(2-\tfrac2\ell\big)\OPT\).
\end{theorem}
\begin{proof}
Let \(S_0=\{r\}\subset S_1\subset\dots\subset S_{k-1}=T\) be the terminal sets after each round and \(\delta_i\) the cost of the
\(i\)-th path. Because the path goes from the current tree, which contains \(S_{i-1}\), to the nearest outside terminal,
\(\delta_i\le\min\{d(s,t):s\in S_{i-1},\,t\notin S_{i-1}\}\), the cheapest edge of \(D\) crossing the cut \((S_{i-1},T\setminus S_{i-1})\).
Let \(M\) be an MST of \(D\). We match round \(i\) to a distinct edge of \(M\) that crosses cut \(i\). For any \(j\) rounds
\(i_1<\dots<i_j\) the cuts split \(T\) into \(j+1\) non-empty layers \(S_{i_1-1},\,S_{i_2-1}\setminus S_{i_1-1},\dots,T\setminus S_{i_j-1}\).
\(M\) is connected, so at least \(j\) of its edges run between different layers, and each of them crosses at least one of the
\(j\) cuts. This is Hall's condition, so the matching exists and \(\sum_i\delta_i\le w(M)\). The clean-up does not increase the
cost, and \cref{thm:kmb} bounds \(w(M)\).
\end{proof}

\subsection{G3: best-gain triples (Zelikovsky)}
For a triple \(z=\{a,b,c\}\subseteq T\) let \(v(z)=\arg\min_v d(a,v)+d(b,v)+d(c,v)\) be its best centre and \(d(z)\) that
minimum (the cheapest star joining \(a,b,c\)). Let \(F\) be the current distance network and \(F[z]\) the network in which
\(a,b\) and \(b,c\) are joined by zero-cost edges. The \emph{win} of \(z\) is
\[\mathrm{win}_F(z)=w(\mathrm{MST}(F))-w(\mathrm{MST}(F[z]))-d(z).\]

\begin{algorithm}[H]
Dijkstra from every terminal; \(F\gets D\); \(Z\gets\emptyset\)\;
\Repeat{\(\max_z\mathrm{win}_F(z)\le0\)}{
  \(M\gets\) Prim on \(F\); compute \(s_{xy}\) = heaviest edge on the \(M\)-path between \(x,y\), for all pairs\;
  \(z^\ast\gets\arg\max_z\mathrm{win}_F(z)\), using \cref{lem:gain} and lazy evaluation of \(d(z)\)\;
  \lIf{\(\mathrm{win}_F(z^\ast)>0\)}{\(F\gets F[z^\ast]\); \(Z\gets Z\cup\{z^\ast\}\)}
}
\(H\gets\) stars of all \(z\in Z\) \(\cup\) shortest paths of the non-zero edges of \(\mathrm{MST}(F)\)\;
\Return \textsc{Clean-Up}\((H)\)\;
\caption{G3 (\texttt{src/greedy3\_zelikovsky.cpp}).}
\end{algorithm}

\begin{lemma}\label{lem:gain}
Let \(M\) be an MST of \(F\), and let \(s_{ab},s_{bc},s_{ac}\) be the heaviest edge weights on the \(M\)-paths between the three
pairs. Then \(w(\mathrm{MST}(F))-w(\mathrm{MST}(F[z]))=\max(s_{ab},s_{bc},s_{ac})+\min(s_{ab},s_{bc},s_{ac})\).
\end{lemma}
\begin{proof}
The three \(M\)-paths form a subtree shaped like a Y with centre \(o\) (one leg may be empty). Let \(\mu_a,\mu_b,\mu_c\) be the
heaviest weights on the legs from \(o\) to \(a,b,c\) (0 for an empty leg). Joining \(a,b,c\) by zero edges is the same as
contracting them into one node, which turns the Y into two independent cycles through \(o\). An MST of \(F[z]\) is obtained from
\(M\) plus the zero edges by deleting one edge on each cycle, and the two deleted edges must lie on two different legs (deleting
two edges of one leg disconnects the part between them). The best choice deletes the heaviest edge of each of the two legs with the
largest maxima, saving the sum of the two largest of \(\mu_a,\mu_b,\mu_c\). Since \(s_{ab}=\max(\mu_a,\mu_b)\) and similarly for the
other pairs, if \(\mu_a\ge\mu_b\ge\mu_c\) the three values are \(\mu_a,\mu_b,\mu_a\), and max plus min is \(\mu_a+\mu_b\). The other
orders are symmetric.
\end{proof}

\paragraph{Running time.} \(k\) Dijkstra runs cost \(O(k(m+n\log n))\). Each round costs \(O(k^2)\) for Prim on the heap and for the
\(s\) matrix (one traversal of \(M\) per terminal), \(O(k^3)\) to scan the triples with \cref{lem:gain}, and \(O(n)\) per centre that
has to be computed. Centres are cached across rounds, so each of the \(\binom{k}{3}\) triples has its centre computed at most once,
\(O(k^3n)\) in all. Every accepted triple merges three components into one, so there are at most \((k-1)/2\) rounds, and the total
is \(O(k(m+n\log n)+k^4+k^3n)\). In practice few centres are computed: a triple is skipped when its MST saving alone cannot beat the
best win so far, or when it cannot beat it even with the lower bound \(d(z)\ge(D_{ab}+D_{bc}+D_{ac})/2\) (each leg of the star lies on two of
the three pairwise paths). On \texttt{b18} (\(k=50\)) the code scans \val{g3.b18.scanned} triples and computes \val{g3.b18.centres} centres.

\begin{theorem}[Zelikovsky 1993]\label{thm:zel}
G3 returns a tree of cost at most \(w(\mathrm{MST}(D))\le(2-2/\ell)\OPT\), and at most \(\tfrac{11}{6}\OPT\).
\end{theorem}
\begin{proof}[Proof of the first bound and outline of the second]
Write \(C(F,Z)=w(\mathrm{MST}(F))+\sum_{z\in Z}d(z)\). At the start \(C=w(\mathrm{MST}(D))\), and each accepted triple lowers \(C\)
by its win, which is positive. The final \(H\) costs at most \(C\) and \cref{lem:cleanup} keeps it there, so the output costs at most
\(w(\mathrm{MST}(D))\). For \(11/6\), Zelikovsky decomposes an optimal tree into full components and shows that a best 3-restricted
Steiner tree, \(\mathrm{smt}_3\), costs at most \(\tfrac53\OPT\). A charging argument over the triples of that tree then shows that
the greedy choice of the best-win triple gives
\(\mathrm{cost}\le\tfrac12\big(w(\mathrm{MST}(D))+\mathrm{smt}_3\big)\le\tfrac12\big(2+\tfrac53\big)\OPT=\tfrac{11}{6}\OPT\).
The full proof is in \cite{zel93}; we do not reproduce it.
\end{proof}
\noindent The best known polynomial ratio is \(\ln4+\varepsilon\approx1.39\) \cite{bgrs}; it needs an LP solver, which the
assignment rules out.

\subsection{Counterexamples}\label{sec:counter}
\paragraph{A tight family for G1 and G2.}
Take a Steiner vertex \(c\) joined to \(\ell\) terminals by edges of weight 10, and join every pair of terminals by an edge of
weight 19 (\cref{fig:star}). The star costs \(10\ell\) and is optimal. Every terminal pair is at distance \(19<20\), so \(D\)
is the uniform 19-clique; G1 and G2 both build a path of \(\ell-1\) direct edges, and the clean-up keeps it because \(c\) is never
reached. Their cost is \(19(\ell-1)\), a ratio of \(1.9(1-1/\ell)\), which approaches the bound \(2-2/\ell\) of
\cref{thm:kmb,thm:sph} as \(19/10\) is replaced by \((2w-1)/w\) with larger \(w\). G3 is optimal on the whole family: every triple has
win \(2\cdot19-30=8>0\), so G3 contracts triples through \(c\), \(c\) enters \(H\), and the clean-up's MST then finds the star.
\Cref{tab:star} gives the measured costs; they match the formulas exactly.

\begin{figure}[H]\centering
\begin{minipage}{0.36\linewidth}\centering
\begin{tikzpicture}[scale=1.05]
\foreach \i in {1,...,6}{\node[circle,draw,fill=accent,text=white,inner sep=1.5pt,font=\scriptsize] (t\i) at ({60*\i}:1.6) {\(t_\i\)};}
\node[circle,draw,inner sep=1.5pt,font=\scriptsize] (c) at (0,0) {\(c\)};
\foreach \i in {1,...,6}{\draw[thick,cG3] (c) -- node[font=\tiny,fill=white,inner sep=0.5pt]{10} (t\i);}
\foreach \i/\j in {1/2,2/3,3/4,4/5,5/6,6/1,1/3,2/4,3/5,4/6,5/1,6/2,1/4,2/5,3/6}{\draw[black!25] (t\i) -- (t\j);}
\foreach \i/\j in {1/2,2/3,3/4,4/5,5/6}{\draw[very thick,cG1,dashed] (t\i) -- (t\j);}
\end{tikzpicture}
\end{minipage}\hfill
\begin{minipage}{0.6\linewidth}
\small\begin{tabular}{@{}rrrrrrr@{}}\toprule
\(\ell\) & OPT & G1 & G2 & G3 & G1/OPT & \(2-2/\ell\)\\\midrule
3 & 30 & 38 & 38 & 30 & 1.267 & 1.333\\
4 & 40 & 57 & 57 & 40 & 1.425 & 1.500\\
5 & 50 & 76 & 76 & 50 & 1.520 & 1.600\\
6 & 60 & 95 & 95 & 60 & 1.583 & 1.667\\
7 & 70 & 114 & 114 & 70 & 1.629 & 1.714\\
8 & 80 & 133 & 133 & 80 & 1.663 & 1.750\\
9 & 90 & 152 & 152 & 90 & 1.689 & 1.778\\
10 & 100 & 171 & 171 & 100 & 1.710 & 1.800\\
11 & 110 & 190 & 190 & 110 & 1.727 & 1.818\\
12 & 120 & 209 & 209 & 120 & 1.742 & 1.833\\
\bottomrule\end{tabular}
\captionof{table}{Measured costs on the star family.}\label{tab:star}

\end{minipage}
\caption{Left: the gadget for \(\ell=6\). Grey edges have weight 19. The dashed path is what G1 and G2 return (cost 95); the star
through \(c\) is optimal (cost 60) and is what G3 returns. Right: measured costs for \(\ell=3,\dots,12\).}\label{fig:star}
\end{figure}

\paragraph{Instances found by search.}
For G3 the star family is no counterexample, and the theorem gives no tight family that is easy to write down. So
\texttt{scripts/gen\_instances.py} searches for one. It starts from small random graphs (7 to 9 vertices, 4 or 5 terminals) and
hill-climbs on edge weights, edges and the terminal set to maximise \(\mathrm{cost}(\mathrm{G}x)/\OPT\), with \(\OPT\) from the
dynamic program, accepting ties so it can cross plateaus. It then deletes edges as long as the ratio does not drop, so that the
graphs stay small enough to draw. Six restarts per strategy, all seeded. \Cref{fig:gadgets} shows the worst instance found
for each strategy. The search found nothing above 1.4 on graphs this small, far from the proven bounds: those bounds need many leaves,
and a 4- or 5-terminal tree has few.

\begin{figure}[!htbp]\centering
\begin{subfigure}{0.32\linewidth}\centering\begin{tikzpicture}[scale=1.25]
\node[circle,draw,fill=white,inner sep=1pt,minimum size=4.6mm,font=\scriptsize] (v0) at (90.0:1.35) {1};
\node[circle,draw,fill=accent,text=white,inner sep=1pt,minimum size=4.6mm,font=\scriptsize] (v1) at (141.4:1.35) {2};
\node[circle,draw,fill=white,inner sep=1pt,minimum size=4.6mm,font=\scriptsize] (v2) at (192.9:1.35) {3};
\node[circle,draw,fill=accent,text=white,inner sep=1pt,minimum size=4.6mm,font=\scriptsize] (v3) at (244.3:1.35) {4};
\node[circle,draw,fill=accent,text=white,inner sep=1pt,minimum size=4.6mm,font=\scriptsize] (v4) at (295.7:1.35) {5};
\node[circle,draw,fill=accent,text=white,inner sep=1pt,minimum size=4.6mm,font=\scriptsize] (v5) at (347.1:1.35) {6};
\node[circle,draw,fill=white,inner sep=1pt,minimum size=4.6mm,font=\scriptsize] (v6) at (398.6:1.35) {7};
\draw[black!30] (v0) -- node[font=\tiny,fill=white,inner sep=0.4pt]{12} (v1);
\draw[cG1,thick,dashed] (v0) -- (v1);
\draw[line width=2.6pt,cG3!45] (v0) -- node[font=\tiny,fill=white,inner sep=0.4pt]{4} (v2);
\draw[line width=2.6pt,cG3!45] (v0) -- node[font=\tiny,fill=white,inner sep=0.4pt]{3} (v3);
\draw[cG1,thick,dashed] (v0) -- (v3);
\draw[line width=2.6pt,cG3!45] (v2) -- node[font=\tiny,fill=white,inner sep=0.4pt]{10} (v4);
\draw[line width=2.6pt,cG3!45] (v1) -- node[font=\tiny,fill=white,inner sep=0.4pt]{5} (v6);
\draw[cG1,thick,dashed] (v1) -- (v6);
\draw[black!30] (v4) -- node[font=\tiny,fill=white,inner sep=0.4pt]{17} (v5);
\draw[cG1,thick,dashed] (v4) -- (v5);
\draw[line width=2.6pt,cG3!45] (v2) -- node[font=\tiny,fill=white,inner sep=0.4pt]{3} (v6);
\draw[line width=2.6pt,cG3!45] (v5) -- node[font=\tiny,fill=white,inner sep=0.4pt]{5} (v6);
\draw[cG1,thick,dashed] (v5) -- (v6);
\end{tikzpicture}
\caption{G1: cost 42, OPT 30, ratio 1.400.}\end{subfigure}\hfill
\begin{subfigure}{0.32\linewidth}\centering\begin{tikzpicture}[scale=1.25]
\node[circle,draw,fill=accent,text=white,inner sep=1pt,minimum size=4.6mm,font=\scriptsize] (v0) at (90.0:1.35) {1};
\node[circle,draw,fill=white,inner sep=1pt,minimum size=4.6mm,font=\scriptsize] (v1) at (135.0:1.35) {2};
\node[circle,draw,fill=accent,text=white,inner sep=1pt,minimum size=4.6mm,font=\scriptsize] (v2) at (180.0:1.35) {3};
\node[circle,draw,fill=white,inner sep=1pt,minimum size=4.6mm,font=\scriptsize] (v3) at (225.0:1.35) {4};
\node[circle,draw,fill=white,inner sep=1pt,minimum size=4.6mm,font=\scriptsize] (v4) at (270.0:1.35) {5};
\node[circle,draw,fill=accent,text=white,inner sep=1pt,minimum size=4.6mm,font=\scriptsize] (v5) at (315.0:1.35) {6};
\node[circle,draw,fill=accent,text=white,inner sep=1pt,minimum size=4.6mm,font=\scriptsize] (v6) at (360.0:1.35) {7};
\node[circle,draw,fill=accent,text=white,inner sep=1pt,minimum size=4.6mm,font=\scriptsize] (v7) at (405.0:1.35) {8};
\draw[line width=2.6pt,cG3!45] (v3) -- node[font=\tiny,fill=white,inner sep=0.4pt]{10} (v6);
\draw[black!30] (v4) -- node[font=\tiny,fill=white,inner sep=0.4pt]{18} (v6);
\draw[cG1,thick,dashed] (v4) -- (v6);
\draw[line width=2.6pt,cG3!45] (v6) -- node[font=\tiny,fill=white,inner sep=0.4pt]{1} (v7);
\draw[cG1,thick,dashed] (v6) -- (v7);
\draw[line width=2.6pt,cG3!45] (v2) -- node[font=\tiny,fill=white,inner sep=0.4pt]{18} (v3);
\draw[black!30] (v2) -- node[font=\tiny,fill=white,inner sep=0.4pt]{9} (v4);
\draw[cG1,thick,dashed] (v2) -- (v4);
\draw[line width=2.6pt,cG3!45] (v3) -- node[font=\tiny,fill=white,inner sep=0.4pt]{10} (v5);
\draw[line width=2.6pt,cG3!45] (v0) -- node[font=\tiny,fill=white,inner sep=0.4pt]{10} (v3);
\draw[black!30] (v5) -- node[font=\tiny,fill=white,inner sep=0.4pt]{20} (v6);
\draw[cG1,thick,dashed] (v5) -- (v6);
\draw[black!30] (v0) -- node[font=\tiny,fill=white,inner sep=0.4pt]{20} (v7);
\draw[cG1,thick,dashed] (v0) -- (v7);
\draw[black!30] (v1) -- node[font=\tiny,fill=white,inner sep=0.4pt]{6} (v7);
\end{tikzpicture}
\caption{G2: cost 68, OPT 49, ratio 1.388.}\end{subfigure}\hfill
\begin{subfigure}{0.32\linewidth}\centering\begin{tikzpicture}[scale=1.25]
\node[circle,draw,fill=accent,text=white,inner sep=1pt,minimum size=4.6mm,font=\scriptsize] (v0) at (90.0:1.35) {1};
\node[circle,draw,fill=white,inner sep=1pt,minimum size=4.6mm,font=\scriptsize] (v1) at (130.0:1.35) {2};
\node[circle,draw,fill=accent,text=white,inner sep=1pt,minimum size=4.6mm,font=\scriptsize] (v2) at (170.0:1.35) {3};
\node[circle,draw,fill=accent,text=white,inner sep=1pt,minimum size=4.6mm,font=\scriptsize] (v3) at (210.0:1.35) {4};
\node[circle,draw,fill=accent,text=white,inner sep=1pt,minimum size=4.6mm,font=\scriptsize] (v4) at (250.0:1.35) {5};
\node[circle,draw,fill=white,inner sep=1pt,minimum size=4.6mm,font=\scriptsize] (v5) at (290.0:1.35) {6};
\node[circle,draw,fill=white,inner sep=1pt,minimum size=4.6mm,font=\scriptsize] (v6) at (330.0:1.35) {7};
\node[circle,draw,fill=accent,text=white,inner sep=1pt,minimum size=4.6mm,font=\scriptsize] (v7) at (370.0:1.35) {8};
\node[circle,draw,fill=white,inner sep=1pt,minimum size=4.6mm,font=\scriptsize] (v8) at (410.0:1.35) {9};
\draw[line width=2.6pt,cG3!45] (v4) -- node[font=\tiny,fill=white,inner sep=0.4pt]{1} (v7);
\draw[cG1,thick,dashed] (v4) -- (v7);
\draw[black!30] (v1) -- node[font=\tiny,fill=white,inner sep=0.4pt]{11} (v7);
\draw[cG1,thick,dashed] (v1) -- (v7);
\draw[line width=2.6pt,cG3!45] (v5) -- node[font=\tiny,fill=white,inner sep=0.4pt]{1} (v6);
\draw[cG1,thick,dashed] (v5) -- (v6);
\draw[line width=2.6pt,cG3!45] (v3) -- node[font=\tiny,fill=white,inner sep=0.4pt]{14} (v5);
\draw[line width=2.6pt,cG3!45] (v0) -- node[font=\tiny,fill=white,inner sep=0.4pt]{11} (v6);
\draw[cG1,thick,dashed] (v0) -- (v6);
\draw[line width=2.6pt,cG3!45] (v6) -- node[font=\tiny,fill=white,inner sep=0.4pt]{10} (v7);
\draw[cG1,thick,dashed] (v6) -- (v7);
\draw[black!30] (v1) -- node[font=\tiny,fill=white,inner sep=0.4pt]{13} (v3);
\draw[cG1,thick,dashed] (v1) -- (v3);
\draw[line width=2.6pt,cG3!45] (v2) -- node[font=\tiny,fill=white,inner sep=0.4pt]{11} (v5);
\draw[cG1,thick,dashed] (v2) -- (v5);
\draw[black!30] (v0) -- node[font=\tiny,fill=white,inner sep=0.4pt]{13} (v8);
\end{tikzpicture}
\caption{G3: cost 58, OPT 48, ratio 1.208.}\end{subfigure}
\caption{Worst instances found by the search, one per strategy. Filled vertices are terminals. The wide light band marks an optimal tree (from the DP); the dashed line is the tree the greedy strategy returns. Numbers on edges are weights.}\label{fig:gadgets}\end{figure}


\subsection{Optimality gaps on the benchmark}\label{sec:gaps}
\Cref{tab:gaps} compares the three strategies on every instance with a known optimum (published for SteinLib, from the dynamic
program for the generated sets). The gap is \((\mathrm{cost}-\OPT)/\OPT\).

\begin{table}[!htbp]\centering\small
\begin{tabular}{@{}lrrrrrrrrrrrrr@{}}\toprule
 & & \multicolumn{4}{c}{G1 KMB} & \multicolumn{4}{c}{G2 SPH (best root)} & \multicolumn{4}{c}{G3 Zelikovsky}\\
\cmidrule(lr){3-6}\cmidrule(lr){7-10}\cmidrule(l){11-14}
Set & \# & mean & max & opt & ms & mean & max & opt & ms & mean & max & opt & ms\\\midrule
B & 18 & 2.38 & 6.30 & 7 & 1.0 & \textbf{0.18} & 1.64 & 15 & 11.3 & 0.65 & 4.72 & 12 & 1.5\\
C & 20 & 5.36 & 15.62 & 1 & 26.1 & \textbf{1.99} & 9.09 & 5 & 1703.0 & 2.33 & 9.09 & 3 & 153.9\\
P6E & 15 & 2.77 & 8.96 & 2 & 1.5 & 0.36 & 3.49 & 10 & 39.1 & \textbf{0.09} & 1.27 & 13 & 1.9\\
I080 & 100 & 18.51 & 41.03 & 1 & 0.6 & 10.73 & 41.03 & 13 & 6.5 & \textbf{8.14} & 20.48 & 3 & 0.9\\
Grid & 27 & 2.89 & 9.41 & 6 & 0.5 & \textbf{0.54} & 2.68 & 17 & 2.4 & 0.57 & 2.72 & 17 & 0.6\\
Random & 6 & 7.95 & 14.27 & 0 & 4.5 & \textbf{1.03} & 2.04 & 1 & 22.1 & 2.42 & 5.30 & 0 & 4.6\\
Adversarial & 13 & 50.48 & 74.17 & 1 & 0.1 & 48.78 & 74.17 & 2 & 0.2 & \textbf{1.60} & 20.83 & 12 & 0.1\\
\midrule All & 199 & 14.19 & 74.17 & 18 & 3.3 & 8.93 & 74.17 & 63 & 179.4 & 4.65 & 20.83 & 60 & 16.4\\
\bottomrule\end{tabular}
\caption{Gap to the optimum in percent (mean and maximum), number of instances solved optimally, and mean time in milliseconds. Bold: smallest mean gap in the row.}\label{tab:gaps}
\end{table}


\noindent First, all three are far better on the benchmarks than their worst case: the largest gap on a SteinLib instance is
\val{gap.max.steinlib}\%, nowhere near a ratio of \(11/6\) or \(2\); the gaps of 50--74\% appear only on the adversarial set, which was built to
cause them. Second, no strategy wins everywhere. G3 has the smallest mean gap overall (\val{g3.meangap}\%) because it is the only one that
stays close to the optimum on the adversarial instances and on I080, whose incidence weights punish pairwise path joins. G2 has the
smallest gap on B, C, grids and random graphs. Third, G2 pays for that with time: it runs the heuristic once per terminal, which makes it
the slowest greedy strategy by an order of magnitude (\val{time.C.G2}\,s per C instance on average, against \val{time.C.G3}\,s for G3 and \val{time.C.G1}\,s for G1). With a single root
it is much weaker, and its worst root is more expensive than G1 on \val{sph.worse_than_g1} of \val{gap.count} instances
(\cref{tab:sphroots}). The choice of root is part of the greedy decision, and a bad one costs as much as the difference between G1 and G2.

\begin{table}[H]\centering\small
\begin{tabular}{@{}lrrrrr@{}}\toprule
Set & G1 & G2, first terminal & G2, worst root & G2, best root & worst \(>\) G1\\\midrule
B & 2.38 & 1.57 & 2.65 & 0.18 & 5/18\\
C & 5.36 & 2.97 & 4.57 & 1.99 & 7/20\\
P6E & 2.77 & 1.17 & 2.53 & 0.36 & 5/15\\
I080 & 18.51 & 14.62 & 17.28 & 10.73 & 22/100\\
Grid & 2.89 & 1.37 & 2.80 & 0.54 & 6/27\\
Random & 7.95 & 3.23 & 4.10 & 1.03 & 2/6\\
Adversarial & 50.48 & 48.78 & 50.38 & 48.78 & 1/13\\
\bottomrule\end{tabular}
\caption{Mean gap (\%) of G2 depending on the root terminal, next to G1. The last column counts instances where the worst root gives a more expensive tree than G1.}\label{tab:sphroots}
\end{table}

